\section{Agents and multi-agent systems}
\label{sec:agents-and-mas}
\slidesrc{Intro to agents \& MAS: 1--66}

Part~I defined AI as the study of \emph{intelligent agents}. This chapter makes the word
\emph{agent} precise, defines \emph{multi-agent systems} (MAS), explains the three ways in
which MAS are used (to \emph{design} software, to \emph{implement} it, to \emph{simulate}
the world), and tells the history of the idea, from AI planning and objects to
intentional systems, modal logic, the Contract Net Protocol and actors.

\subsection{What is an agent?}
\label{sec:what-is-an-agent}
\slidesrc{Intro to agents \& MAS: 5--11}

There is no single definition of agent, but two families of definitions are standard.

\subsubsection{The weak definition}
\slidesrc{Intro to agents \& MAS: 5--8}

\begin{nbdef}{Agent, weak notion (Jennings, Sycara, Wooldridge, 1998)}
An agent is a hardware or software system that is
\begin{itemize}[nosep]
  \item \textbf{situated},
  \item \textbf{autonomous},
  \item \textbf{flexible}, which means:
        \begin{itemize}[nosep]
          \item \textbf{reactive},
          \item \textbf{proactive},
          \item \textbf{social}.
        \end{itemize}
\end{itemize}
\end{nbdef}

The slides only list the words; here is what each one means (from the 1998 paper by
N.~Jennings, K.~Sycara and M.~Wooldridge, whose photos are on slide 8).

\begin{center}
\small
\renewcommand{\arraystretch}{1.2}
\begin{tabularx}{\linewidth}{@{}l >{\raggedright\arraybackslash}X >{\raggedright\arraybackslash}X@{}}
\toprule
\textbf{Property} & \textbf{Meaning} & \textbf{Robot vacuum cleaner} \\
\midrule
situated & it lives \emph{in} an environment: it receives sensory input from it and
  acts to change it & senses walls, dirt and stairs; moves and sucks \\
autonomous & it acts without the direct intervention of humans or other agents, and has
  control over its own actions and internal state & nobody drives it: it decides where to
  go next \\
reactive & it perceives the environment and responds \emph{in a timely fashion} to the
  changes that occur in it & a chair is moved: it turns to avoid it \\
proactive & it does not simply act in response to the environment: it exhibits
  \emph{goal-directed} behaviour and takes the initiative & when the battery is low, it
  goes back to the dock on its own initiative \\
social & it interacts with other agents (and humans) to complete its tasks and to help
  others & coordinates with the mop robot so they don't clean the same room \\
\bottomrule
\end{tabularx}
\captionof{table}{The properties of the weak definition, explained with one everyday agent.}
\end{center}

\begin{nbintuition}[The hard part is the balance]
Being only reactive is easy (a thermostat), being only proactive is easy (a program that
blindly executes a plan). The difficulty, and the reason why agent architectures exist
(Section~\ref{sec:agent-architectures}), is to be \emph{both}: pursue goals over time, while
still reacting quickly when the world changes and the plan no longer makes sense.
\end{nbintuition}

\begin{nbexample}{is it an agent?}
\notinslides
\small
\begin{tabularx}{\linewidth}{@{}l >{\raggedright\arraybackslash}X@{}}
\toprule
\textbf{System} & \textbf{Verdict} \\
\midrule
a pocket calculator & no: not situated in an environment it perceives, no autonomy, no goals of its own \\
a thermostat & a degenerate agent: situated, autonomous, reactive, but not proactive nor social \\
a spam filter & situated (in the mail stream), autonomous and reactive; weakly proactive \\
a negotiating bot on an e-commerce platform & yes: it pursues the goal of buying cheaply (proactive), reacts to offers (reactive) and negotiates with sellers' bots (social) \\
\bottomrule
\end{tabularx}
\end{nbexample}

\subsubsection{The strong definition}
\slidesrc{Intro to agents \& MAS: 9--11}

Besides the properties of the weak definition, an agent may be conceptualised following an
\textbf{anthropomorphic} approach (the \textbf{strong} definition): we describe and build it
using notions that we normally apply to humans.

\begin{nbdef}{Agent, strong notion}
\begin{itemize}[nosep]
  \item \textbf{Mentalistic notions} (Y.~Shoham, \emph{Artificial Intelligence} 1993; A.~S.~Rao
        and M.~P.~Georgeff, KR\&R 1992): the agent has \emph{beliefs}, \emph{desires},
        \emph{intentions}, \emph{commitments}, \ldots
  \item \textbf{Emotional notions} (J.~Bates, \emph{Communications of the ACM} 1994):
        friendliness, trust, distrust; agents must be \emph{believable}.
\end{itemize}
\end{nbdef}

\begin{nbexample}{a strong agent: a personal travel assistant}
\notinslides
\begin{itemize}[nosep]
  \item \emph{belief}: ``the 8:15 train to Milan is cancelled'';
  \item \emph{desire}: ``my user should be in Milan before 11:00'';
  \item \emph{intention}: ``book the 8:45 train'' (the desire it committed to, with a plan);
  \item \emph{commitment}: having promised the user to handle the trip, it keeps trying
        alternatives (a bus, a car-sharing ride) until the trip is booked or impossible.
\end{itemize}
Beliefs, desires and intentions are exactly the three components of the \textbf{BDI}
architecture (Section~\ref{sec:bdi-agentspeak}).
\end{nbexample}

\begin{nbwhy}[Believable agents]
Emotional notions matter wherever agents interact with people: virtual characters in video
games and training simulations, conversational assistants, social robots in elderly care.
A character that reacts with plausible emotions is ``believable'', and users trust it
(or distrust it) accordingly.
\end{nbwhy}

\subsection{What is a multi-agent system?}
\label{sec:what-is-a-mas}
\slidesrc{Intro to agents \& MAS: 12--17}

A single agent is rarely enough: interesting problems require several agents, each with
its own goals and abilities, that interact. Such a system is a \textbf{multi-agent system}
(MAS). Slide 12 shows the classic picture, redrawn in Figure~\ref{fig:mas-structure}.

\begin{center}
\begin{tikzpicture}[font=\small, >=Stealth,
  ag/.style={ellipse, draw=nbgreen!60!black, fill=nbgreen!35, thick, minimum width=1.4cm, minimum height=0.75cm, font=\small\bfseries}]
  % environment
  \draw[thick, nbgray, fill=nbgray!15] (5.2,-2.6) ellipse (5.0cm and 0.9cm);
  \node[font=\small\bfseries, nbgray!70!black] at (5.2,-3.85) {Environment};
  \fill[nbblue!25, draw=nbblue] (2.4,-2.55) ellipse (0.9cm and 0.3cm);
  \fill[nbblue!25, draw=nbblue] (4.1,-2.55) ellipse (0.9cm and 0.3cm);
  \fill[nbblue!25, draw=nbblue] (7.6,-2.45) ellipse (0.8cm and 0.3cm);
  % organizations
  \draw[thick, nbgray] (2.9,0.6) ellipse (2.7cm and 1.35cm);
  \draw[thick, nbgray] (7.1,0.75) ellipse (2.9cm and 1.45cm);
  \node[font=\small\bfseries] at (0.4,2.1) {Organisation};
  % agents
  \node[ag] (a1) at (1.1,0.2) {Agent};
  \node[ag] (a2) at (3.0,1.2) {Agent};
  \node[ag] (a3) at (5.0,0.6) {Agent};
  \node[ag] (a4) at (6.9,1.6) {Agent};
  \node[ag] (a5) at (9.0,1.0) {Agent};
  \node[ag] (a6) at (7.4,0.0) {Agent};
  \draw[<->, thick] (a1) -- (a2) node[midway, below right=-2pt, font=\scriptsize, align=left]{inter-agent\\communication};
  \draw[<->, thick] (a2) -- (a3);
  \draw[<->, thick] (a3) -- (a4);
  \draw[<->, thick] (a4) -- (a5);
  \draw[<->, thick] (a4) -- (a6);
  \draw[<->, thick] (a3) -- (a6);
  % access
  \draw[<->, ultra thick, nborange] (1.4,-0.3) -- (2.3,-2.2);
  \draw[<->, ultra thick, nborange] (4.8,0.15) -- (4.2,-2.2);
  \draw[<->, ultra thick, nborange] (7.5,-0.45) -- (7.6,-2.1);
  \node[font=\scriptsize, align=left] at (9.2,-1.4) {access to the\\environment\\(sphere of influence)};
\end{tikzpicture}
\captionof{figure}{Structure of a MAS (slide 12). Agents (green) communicate with each other,
are grouped in organisations (grey ellipses) and each one can perceive and act on its own
portion of the shared environment (blue regions, orange arrows). Spheres of influence may
overlap: two agents may act on the same part of the environment.}
\label{fig:mas-structure}
\end{center}

\begin{nbdef}{Characteristics of a MAS (slides 14--17)}
In a multi-agent system:
\begin{enumerate}[nosep]
  \item each agent has \textbf{incomplete information}, or capabilities for solving the
        problem, thus each agent has a \textbf{limited viewpoint};
  \item there is \textbf{no global system control};
  \item data is \textbf{decentralised};
  \item computation is \textbf{asynchronous}.
\end{enumerate}
\end{nbdef}

\begin{nbexample}{a smart-city ride-sharing MAS}
\notinslides
Passenger agents (one per user app), driver agents (one per car) and a pricing agent.
\begin{enumerate}[nosep]
  \item \emph{Limited viewpoint}: a driver agent knows its own position and route, not the
        positions of the other cars or the requests of all passengers.
  \item \emph{No global control}: no central program decides every move; drivers accept or
        refuse rides.
  \item \emph{Decentralised data}: each phone stores its own user's preferences and history.
  \item \emph{Asynchronous computation}: requests and answers arrive at any time; nobody
        waits for a global ``clock tick''.
\end{enumerate}
The same four features describe e-commerce marketplaces, power grids with home batteries,
fleets of warehouse robots and swarms of drones.
\end{nbexample}

\begin{nbintuition}[Why ``no global control'' is a feature, not a bug]
A central controller is a single point of failure and a bottleneck, and in an \emph{open}
system (agents of different owners, joining and leaving at any time) it simply cannot exist:
nobody owns all the agents. MAS replace central control by \emph{interaction}: negotiation,
coordination, competition. This is why agent \emph{communication}
(Section~\ref{sec:agent-communication}) is a central topic of the course.
\end{nbintuition}

\subsection{Three ways of using MAS}
\label{sec:uses-of-mas}
\slidesrc{Intro to agents \& MAS: 18--41}

Agents and MAS can enter a software project at three different levels:

\begin{center}
\begin{tikzpicture}[font=\small, >=Stealth,
  b/.style={draw, rounded corners, thick, align=center, text width=4.0cm, minimum height=2.0cm}]
  \node[b, draw=nbblue, fill=nbblue!7] (se) at (0,0) {\textbf{1. SE metaphor}\\analysis and design\\think of the system as agents\\[2pt]\footnotesize PDT, INGENIAS};
  \node[b, draw=nbpurple, fill=nbpurple!7] (dev) at (4.7,0) {\textbf{2. Development}\\implementation\\the code \emph{is} agents\\[2pt]\footnotesize JADE, Jason};
  \node[b, draw=nbgreen, fill=nbgreen!7] (sim) at (9.4,0) {\textbf{3. Simulation}\\model a real system\\as interacting agents\\[2pt]\footnotesize NetLogo};
\end{tikzpicture}
\captionof{figure}{The three uses of MAS (slides 18--41). They can be combined: e.g.\ design with
agents and implement with mainstream technologies.}
\end{center}

\subsubsection{MAS as a software engineering metaphor}
\slidesrc{Intro to agents \& MAS: 18--21}

\begin{nbfact}{Zambonelli and Omicini (JAAMAS, 2004)}
``\ldots today's software engineering approaches are increasingly adopting abstractions
approaching that of agent-based computing. This trend can be better understood by
recognising that \emph{the vast majority of modern distributed systems scenarios are
intrinsically prone to be developed in terms of MASs}, and that modern distributed systems
are already de facto MASs, i.e., they are indeed composed of autonomous, situated, and
social components.''
\end{nbfact}

In other words: even if you never write a line of ``agent code'', thinking in terms of
agents (who are the autonomous components? what are their goals? how do they interact?)
helps in the \emph{analysis} and \emph{design} of distributed systems. Tools exist to support
this: \textbf{PDT} (Prometheus Design Tool, RMIT University) and \textbf{INGENIAS}, which
let you draw agents, goals, messages and data stores as diagrams. Methodologies are the
topic of Section~\ref{sec:engineering-mas}.

\subsubsection{Agent-based development}
\slidesrc{Intro to agents \& MAS: 22--36}

\begin{nbdef}{Agent-based development (slides 22--23)}
Agents are employed as the technology for \textbf{implementing the software application
itself}. This choice is usually made when
\begin{itemize}[nosep]
  \item the agent metaphor was already applied in the earlier engineering steps (analysis,
        design);
  \item the intelligence of the \emph{individual} agent is more relevant than the
        ``emerging intelligence'' of the overall system;
  \item complex patterns of interaction may be required;
  \item the application is geographically \textbf{distributed}, \textbf{decentralised}, or
        \textbf{open}.
\end{itemize}
\end{nbdef}

Two agent platforms are shown: \textbf{JADE} (Java Agent DEvelopment framework,
Section~\ref{sec:jade}), whose sniffer tool draws the messages exchanged by agents as a
sequence diagram (slide 24), and \textbf{Jason} (slide 25, an interpreter of AgentSpeak, see
the book \emph{Programming multi-agent systems in AgentSpeak using Jason};
Section~\ref{sec:bdi-agentspeak}). (Slide 25 is labelled ``Jade'' but refers to Jason.)

\begin{nbwarning}[Why industry rarely uses agent platforms (slides 26--27)]
Many agent-based development tools exist, but they are \emph{not commercial projects
supported by large and solid companies}. This is one of the reasons for the limited adoption
of this technology by industry. A pragmatic answer: use technologies that are not
officially tagged as ``agent-oriented'', but are very suitable for implementing software
agents, in the \emph{implementation} stage, and a ``pure'' agent approach for the
\emph{modelling and design} stages.
\end{nbwarning}

The slides then list research works that follow this idea, mapping agents onto mainstream
technologies over the years:

\begin{center}
\small
\begin{tabularx}{\linewidth}{@{}l l X@{}}
\toprule
\textbf{Year} & \textbf{Technology} & \textbf{Work shown on the slide} \\
\midrule
2009 & service-oriented computing & \emph{Agents and Service-Oriented Computing for Autonomic Computing: A Research Agenda} \\
2011 & cloud & \emph{Software Agents as Cloud Computing Services}; \emph{Cloud Computing and Software Agents: Towards Cloud Intelligent Services} \\
2013 & web services & JADE Web Services Integration Gateway (WSIG) \\
2013 & Internet of Things & \emph{From Internet of Things to Internet of Agents} \\
2015 & Internet of Things & \emph{Agents of Things (AoT): an intelligent operational concept of the IoT} \\
2018 & IoT and the Web & EMAS 2018 keynote: \emph{Autonomous Agents for Flexible Hypermedia Systems} \\
2023 & virtual reality & \emph{VEsNA}: virtual environments via natural language agents, applied to factory automation (Gatti, Mascardi) \\
2024 & extended reality & XRIA 2024: first workshop on ``eXtended Reality \& Intelligent Agents'' at ECAI \\
\bottomrule
\end{tabularx}
\captionof{table}{Agents implemented with mainstream technologies (slides 28--36).}
\end{center}

\begin{nbwhy}[Agents everywhere, under other names]
A microservice that watches a queue, decides and calls other services is an agent in
disguise; so is an IoT device that negotiates energy prices, or a non-player character in a
VR training environment that understands spoken commands. The agent \emph{concepts}
(autonomy, goals, interaction) transfer to these technologies even when no agent platform is
used.
\end{nbwhy}

\subsubsection{MAS for simulation}
\slidesrc{Intro to agents \& MAS: 37--41}

\begin{nbfact}{Agent-based modelling and simulation (Macal and North, 2007)}
``Agent-based modeling and simulation (ABMS) is a new approach to modeling systems
comprised of autonomous, interacting agents. ABMS promises to have far-reaching effects on
the way that businesses use computers to support decision-making and researchers use
electronic laboratories to support their research.''

``Some have gone so far as to contend that ABMS is a \emph{third way of doing science}, in
addition to traditional deductive and inductive reasoning. [\ldots] Applications range from
modeling agent behavior in the stock market, supply chains, and consumer markets, to
predicting the spread of epidemics, the threat of bio-warfare, and the factors responsible
for the fall of ancient civilizations.''
\end{nbfact}

The tool shown on slide 41 is \textbf{NetLogo}, a programmable environment where thousands
of simple agents (``turtles'') move on a grid of ``patches'' and the global behaviour is
observed in real time. ABMS is the topic of Section~\ref{sec:abms}.

\begin{nbintuition}[Deduction, induction, simulation]
Deduction: from general laws to particular consequences (logic, Part~II). Induction: from
many observations to general laws (statistics, machine learning). Simulation: write the
\emph{local} rules of each individual, run them, and \emph{observe} what global phenomenon
emerges (a traffic jam, an epidemic wave, a market crash). That is the ``third way''.
\end{nbintuition}

\subsection{A short history of agents and MAS}
\label{sec:agents-history}
\slidesrc{Intro to agents \& MAS: 42--66}

Current interest in autonomous agents did not emerge from a vacuum. The main contributing
disciplines are \textbf{artificial intelligence}, \textbf{object-oriented programming} (and
concurrent object-based systems) and \textbf{human-computer interface design} (slide 43).

\subsubsection{From artificial intelligence}
\slidesrc{Intro to agents \& MAS: 44--47}

\begin{center}
\begin{tikzpicture}[font=\small, >=Stealth,
  ev/.style={anchor=north west, text width=10.2cm, align=left, inner sep=2pt},
  yr/.style={anchor=north east, font=\small\bfseries, inner sep=2pt}]
  \coordinate (e0) at (0.3,0);
  \foreach \i/\year/\txt/\c in {
    1/1961--71/{AI \textbf{planning}: Newell and Simon (1961), Fikes and Nilsson's STRIPS (1971). Until the 1980s, AI studies the components of intelligence (learning, reasoning, vision\ldots) \emph{in isolation}; expert systems are ``disembodied intelligence'': no real-time capabilities, no social ability.}/nbblue,
    2/1987/{Disillusion with \textbf{first-principle planning} (Chapman): building a new plan from scratch for every goal is too slow to be reactive (``calculative rationality'') and undecidable in many cases. Georgeff and Lansky propose the \textbf{BDI} architecture for practical reasoning.}/nbred,
    3/1991/{Brooks: \textbf{behavioural} or \textbf{reactive AI}, intelligence without symbolic representation.}/nborange,
    4/1992--94/{\textbf{Hybrid architectures}, usually stratified in layers: TouringMachines (Ferguson, 1992), InteRRaP (M\"uller and Pischel, 1994).}/nbpurple,
    5/1993/{Shoham, \emph{Agent-oriented programming}: agents are not just a branch of AI, but an autonomous technological and methodological approach. Shoham is considered the father of agent-oriented programming.}/nbgreen}
  {
    \pgfmathtruncatemacro{\prev}{\i-1}
    \node[ev] (e\i) at ([yshift=-4pt]e\prev.south west) {\txt};
    \node[yr, text=\c] at ([xshift=-0.55cm]e\i.north west) {\year};
    \fill[\c] ([xshift=-0.3cm, yshift=-7pt]e\i.north west) circle (3.2pt);
  }
  \begin{scope}[on background layer]
    \draw[very thick, nbblue!50] (0,0) -- (0,0 |- e5.south);
  \end{scope}
\end{tikzpicture}
\captionof{figure}{From AI to agents (slides 44--47). The architectures named here (reactive,
deliberative, hybrid, BDI) are the subject of Section~\ref{sec:agent-architectures}.}
\end{center}

\begin{nbdef}{Practical reasoning agents (slide 46)}
Agents whose architecture is modelled on, or inspired by, a theory of \textbf{practical
reasoning} in humans: the kind of pragmatic reasoning that we use to decide \emph{what to
do} (as opposed to theoretical reasoning, which decides what is true). The most widely known
practical reasoning architecture is \textbf{Belief-Desire-Intention} (BDI).
\end{nbdef}

\subsubsection{From object-oriented programming}
\slidesrc{Intro to agents \& MAS: 48--49}

Objects are computational entities that \emph{encapsulate some state}, are able to
\emph{perform actions} (methods) on this state, and \emph{communicate by message passing}.
This sounds very much like an agent. So, what is the difference?

\begin{center}
\small
\begin{tabularx}{\linewidth}{@{}l >{\raggedright\arraybackslash}X >{\raggedright\arraybackslash}X@{}}
\toprule
 & \textbf{Object} & \textbf{Agent} \\
\midrule
Autonomy & controls its \emph{state} (private fields), but not its \emph{behaviour}: if
  another object calls a public method, the method runs & controls its state \emph{and} its
  behaviour: a request is a message the agent may refuse \\
Motto & \multicolumn{2}{c}{``Objects do it for free; agents do it for money.''} \\
Flexibility & the standard object model says nothing about how to build systems that
  integrate proactive, reactive and social behaviour & reactive, proactive and social
  behaviour are part of the model \\
Thread of control & usually passive: acts only when called & has its own thread of
  control: always running, takes the initiative \\
\bottomrule
\end{tabularx}
\captionof{table}{Agents vs.\ objects (slide 49; last row added).}
\end{center}

\begin{nbexample}{the same request, two worlds}
\notinslides
Object: \code{account.withdraw(100)} \emph{always} executes the method (perhaps throwing
an exception). Agent: Alice's banking agent \emph{receives} the message ``please transfer
100 euro to Bob'', checks its own goals (``keep at least 500 on the account'', ``do not pay
unknown parties'') and may answer ``refuse''. One could argue that we can program this with
objects too; the point (slide 49) is that the object model has nothing to say about it,
while for agents it is the core of the abstraction.
\end{nbexample}

\subsubsection{From human-computer interaction}
\slidesrc{Intro to agents \& MAS: 50}

Instead of \emph{direct manipulation} (the user does everything by clicking), we would like
programs that behave like \textbf{``digital butlers''} (N.~Negroponte, \emph{Being
Digital}, 1995):
\begin{itemize}[nosep]
  \item ``The agent answers the phone, recognizes the callers, disturbs you when
        appropriate, and may even tell a white lie on your behalf.''
  \item ``If you have somebody who knows you well and shares much of your information,
        that person can act on your behalf very effectively. If your secretary falls ill, it
        would make no difference if the temping agency could send you Albert Einstein. This
        issue is not about IQ. It is shared knowledge and the practice of using it in your
        best interests.''
  \item ``Like an army commander sending a scout ahead \ldots\ you will dispatch agents to
        collect information on your behalf.''
\end{itemize}
Thirty years later, personal assistants and AI agents that read your e-mail and book your
travel are trying to fulfil exactly this vision.

\subsubsection{Agents as intentional systems}
\slidesrc{Intro to agents \& MAS: 51--60}

When explaining human activity we often say things like:
\begin{itemize}[nosep]
  \item Janine took her umbrella because she \emph{believed} it was going to rain;
  \item Michael worked hard because he \emph{wanted} to possess a PhD.
\end{itemize}
These statements use a \textbf{folk psychology}: human behaviour is predicted and explained
by attributing attitudes such as believing, wanting, hoping, fearing. These attitudes are
called the \textbf{intentional notions}.

\begin{nbdef}{Intentional system (D.~C.~Dennett, \emph{The Intentional Stance}, 1989)}
An \textbf{intentional system} is a system made up of entities whose behaviour can be
predicted by the method of attributing \emph{beliefs}, \emph{desires} and \emph{rational
acumen}.
\end{nbdef}

\begin{nbintuition}[The intentional stance in practice]
You do not know the code of a chess program, but you predict its moves by saying ``it
\emph{wants} to win and \emph{believes} my queen is unprotected, so it will take it''. That
works much better than reasoning about transistors. Taking the intentional stance towards a
system is useful when it is too complex to explain mechanically, yet behaves rationally.
\end{nbintuition}

If we adhere to the \emph{strong} definition of agents, intentional systems are a suitable
theory for agents (slides 56--60). Representing intentional notions raises delicate
syntactic and semantic questions, and the usual answers are:
\begin{itemize}[nosep]
  \item \textbf{modal logic} languages are suitable for \emph{specifying} agents as
        intentional systems;
  \item languages based on \textbf{computational (modal) logic} are suitable for
        \emph{programming} agents;
  \item axiomatised logic-based languages can undergo \textbf{axiomatic verification};
        other languages can be used for \textbf{model checking}.
\end{itemize}

\begin{nbdef}{Modal logic (slides 61--63)}
Modal logic extends classical logic with a new connective $\Box$ (\textbf{necessity}) and
its derivable counterpart $\Diamond$ (\textbf{possibility}).
\begin{itemize}[nosep]
  \item $\Box p$: $p$ is \emph{necessarily} true, i.e.\ true in every possible scenario
        (``possible world'');
  \item $\Diamond p$: $p$ is \emph{possibly} true, i.e.\ true in at least one possible
        scenario.
\end{itemize}
$\Diamond$ can be defined from $\Box$: \quad $\Diamond p \liff \neg\Box\neg p$.
\end{nbdef}

``$p$ is possible'' = ``it is not necessary that $p$ is false''. It is the same duality as
$\exists$ and $\forall$ (Section~\ref{sec:quantifiers}): $\Box$ is a ``for all worlds'',
$\Diamond$ a ``there exists a world''.

\begin{center}
\small
\begin{tabularx}{\linewidth}{@{}l l >{\raggedright\arraybackslash}X@{}}
\toprule
\textbf{Modal logic} & \textbf{$\Box p$ reads} & \textbf{Example} \\
\midrule
epistemic & the agent knows/believes $p$ & $B_{\mathit{robot}}\,\mathit{door\_open}$: the robot believes the door is open \\
temporal & $p$ always holds (in the future) & $\Box\neg\mathit{collision}$: the drones never collide \\
deontic & $p$ is obligatory & $O(\mathit{pay\_taxes})$; $\Diamond$ becomes ``permitted'' \\
dynamic & after every execution of action $a$, $p$ holds & $[\mathit{open}]\,\mathit{door\_open}$ \\
\bottomrule
\end{tabularx}
\captionof{table}{Kinds of modal logics (slide 64; readings and examples added). Combinations
exist: \textbf{BDI logic} combines beliefs, desires, intentions and time; \textbf{KARO}
combines knowledge, abilities, results and opportunities.}
\end{center}

\subsubsection{Two ancestors of MAS: Contract Net and Actors}
\slidesrc{Intro to agents \& MAS: 65--66}

\begin{nbdef}{Contract Net Protocol (Davis and Smith, 1983)}
Agents can dynamically take two \textbf{roles}: \textbf{manager} or \textbf{contractor}.
\begin{enumerate}[nosep]
  \item Given a task, an agent first determines whether it can break it into subtasks that
        could be performed concurrently.
  \item As \emph{manager}, it uses the protocol to \textbf{announce} the tasks that could
        be transferred, and requests \textbf{bids} from nodes that could perform them.
  \item A node that receives an announcement replies with a \textbf{bid}, indicating how
        well it thinks it can perform the task.
  \item The manager collects the bids and \textbf{awards} the task to the best bidder,
        which becomes a \emph{contractor}.
\end{enumerate}
\end{nbdef}

\begin{center}
\begin{tikzpicture}[font=\small, >=Stealth, x=1.05cm, y=-0.8cm]
  \foreach \x/\n in {0/Manager, 3.6/Contractor A, 6.6/Contractor B, 9.6/Contractor C}{
    \node[draw, rounded corners, thick, fill=nbblue!8, minimum width=2.2cm] at (\x,0) {\n};
    \draw[dashed, nbgray] (\x,0.6) -- (\x,8.6);
  }
  \draw[->, thick, nbblue] (0,1.3) -- node[above, pos=0.2, font=\scriptsize]{announce task $T$} (3.6,1.3);
  \draw[->, thick, nbblue] (0,1.3) -- (6.6,1.6);
  \draw[->, thick, nbblue] (0,1.3) -- (9.6,1.9);
  \draw[->, thick, nborange] (3.6,3.2) -- node[above, font=\scriptsize]{bid: 5 min} (0,3.2);
  \draw[->, thick, nborange] (6.6,4.0) -- node[above, pos=0.7, font=\scriptsize]{bid: 2 min} (0,4.0);
  \draw[->, thick, nborange] (9.6,4.8) -- node[above, pos=0.8, font=\scriptsize]{bid: 9 min} (0,4.8);
  \node[draw, fill=nbgreen!10, font=\scriptsize, align=center] at (0,6.0) {evaluate\\bids};
  \draw[->, thick, nbgreen!70!black] (0,7.2) -- node[above, pos=0.6, font=\scriptsize]{award $T$} (6.6,7.2);
  \draw[->, thick, nbgray] (6.6,8.2) -- node[above, pos=0.5, font=\scriptsize]{result} (0,8.2);
\end{tikzpicture}
\captionof{figure}{A run of the Contract Net Protocol: one announcement, three bids, the task is
awarded to B (the best bid). The same protocol appears in FIPA as an interaction protocol
(Section~\ref{sec:agent-communication}).}
\end{center}

\begin{nbwarning}[Slide 65: who collects the bids?]
The slide says ``\emph{the contractor} collects the bids and awards the task to the best
bidder''. It is the \textbf{manager} who announces the task, collects the bids and awards
the task; the winner then becomes the contractor.
\end{nbwarning}

\begin{nbwhy}
Contract Net is everywhere where work is distributed by auction: a cloud scheduler choosing
the server that can run a job fastest, a logistics platform assigning a delivery to the
nearest courier, a factory controller assigning a part to the least loaded machine.
\end{nbwhy}

\begin{nbdef}{Actors (Hewitt, Bishop, Steiger, IJCAI 1973; Agha, 1986)}
Actors were proposed as universal primitives of concurrent computation: self-contained,
interactive, autonomous components of a computing system that communicate by
\textbf{asynchronous message passing}. The basic actor primitives are:
\begin{itemize}[nosep]
  \item \textbf{create}: create an actor from a behaviour description and a set of
        parameters (possibly including existing actors);
  \item \textbf{send}: send a message to an actor;
  \item \textbf{become}: change an actor's local state.
\end{itemize}
\end{nbdef}

\begin{nbintuition}[Actors today]
Each actor has a mailbox, processes one message at a time, and can only affect others by
sending messages: no shared memory, so no locks. Erlang (used for telecom switches and
WhatsApp's servers) and the Akka framework are built on actors. Agents add to actors
\emph{goals} and \emph{reasoning}; the messaging core is the same, and JADE agents
(Section~\ref{sec:jade}) work exactly like this.
\end{nbintuition}


\begin{nbsummary}[Agents and MAS]
\begin{itemize}[nosep]
  \item \textbf{Weak definition} (Jennings, Sycara, Wooldridge 1998): an agent is a hardware
        or software system that is \textbf{situated}, \textbf{autonomous} and
        \textbf{flexible}, i.e.\ \textbf{reactive}, \textbf{proactive} and \textbf{social}.
  \item \textbf{Strong definition}: anthropomorphic; \emph{mentalistic} notions (beliefs,
        desires, intentions, commitments: Shoham 1993, Rao \& Georgeff 1992) and
        \emph{emotional} notions (believable agents: Bates 1994).
  \item \textbf{MAS}: each agent has incomplete information/capabilities (limited
        viewpoint); no global control; decentralised data; asynchronous computation.
        Agents live in organisations and act on a shared environment.
  \item Three uses: MAS as a \textbf{software-engineering metaphor} (analysis/design: PDT,
        INGENIAS); \textbf{agent-based development} (implementation: JADE, Jason; or
        mainstream tech such as SOC, cloud, IoT, VR/XR); \textbf{agent-based simulation}
        (ABMS, NetLogo).
  \item Roots: AI (planning $\to$ reactive (Brooks) $\to$ hybrid $\to$ BDI practical
        reasoning; Shoham 1993 = agent-oriented programming), OOP, HCI (``digital
        butlers'').
  \item \textbf{Agents vs.\ objects}: objects control their state, agents also control
        their \emph{behaviour} (``objects do it for free, agents do it for money'') and
        integrate reactive, proactive and social behaviour.
  \item \textbf{Intentional stance} (Dennett): predict behaviour by attributing beliefs,
        desires, rationality. Specified with \textbf{modal logics} ($\Box$ necessity,
        $\Diamond$ possibility, $\Diamond p \liff \neg\Box\neg p$).
  \item \textbf{Contract Net} (1983): manager announces task, contractors bid, manager
        awards to the best bid. \textbf{Actors} (1973): create, send, become; asynchronous
        message passing.
\end{itemize}
\end{nbsummary}
